\bookchapter{Summation notation}\label{ch:summation}

\section*{Explanation}
Adding many numbers is so common that mathematicians use a short symbol, the capital Greek letter sigma:
\[
\sum_{i=1}^{n}a_i=a_1+a_2+\dots+a_n.
\]
The letter $i$ is the \emph{index}, the number below $\Sigma$ is where it starts and the number above is where it ends. A sum is the same as a loop in a program:
\begin{verbatim}
total = 0
for i in range(1, 5):      # i = 1, 2, 3, 4
    total += i ** 2
print(total)               # 30
\end{verbatim}

\begin{example}
$\displaystyle\sum_{i=1}^{4}i^2=1+4+9+16=30$.
\end{example}

\textbf{Rules.} A constant factor can be pulled out, and sums can be split:
\[
\sum_{i=1}^n c\,a_i=c\sum_{i=1}^n a_i,\qquad
\sum_{i=1}^n(a_i+b_i)=\sum_{i=1}^n a_i+\sum_{i=1}^n b_i,\qquad
\sum_{i=1}^n c=n\,c.
\]
A useful formula is $\displaystyle\sum_{i=1}^n i=\frac{n(n+1)}{2}$.

\section*{Three sums that appear everywhere in AI}
\begin{itemize}
  \item \textbf{Weighted sum} (what a neuron computes): $z=\displaystyle\sum_{i=1}^n w_ix_i+b$.
  \item \textbf{Mean}: $\bar x=\dfrac1n\displaystyle\sum_{i=1}^n x_i$.
  \item \textbf{Mean squared error} between predictions $\hat y_i$ and true values $y_i$:
  \[
  \mathrm{MSE}=\frac1n\sum_{i=1}^n(y_i-\hat y_i)^2.
  \]
\end{itemize}

\begin{example}[Weighted sum]
With $w=(0.5,-1,2)$ and $x=(4,2,1)$: $z=0.5\cdot4-1\cdot2+2\cdot1=2-2+2=2$.
\end{example}

\begin{example}[Mean squared error]
Five data points $(x_i,y_i)$ are $(1,2),(2,2),(3,5),(4,4),(5,6)$ and the model predicts $\hat y=x$. The errors $y_i-\hat y_i$ are $1,0,2,0,1$ (the red segments in Figure~\ref{fig:b1_9}). Then
\[
\mathrm{MSE}=\frac{1^2+0^2+2^2+0^2+1^2}{5}=\frac65=1.2.
\]
\end{example}

\begin{center}
\begin{tikzpicture}
\begin{axis}[gr, width=8cm,height=6.5cm, xmin=0,xmax=6,ymin=0,ymax=7, xtick={0,...,6}, ytick={0,...,7}, legend pos=north west]
  \addplot[blue,domain=0:6]{x};
  \addlegendentry{model $\hat y=x$}
  \addplot[only marks,mark=*,red,mark size=2.5pt] coordinates {(1,2)(2,2)(3,5)(4,4)(5,6)};
  \addlegendentry{data $(x_i,y_i)$}
  \draw[red!70!black,very thick] (axis cs:1,1) -- (axis cs:1,2);
  \draw[red!70!black,very thick] (axis cs:3,3) -- (axis cs:3,5);
  \draw[red!70!black,very thick] (axis cs:5,5) -- (axis cs:5,6);
  \node[right,font=\scriptsize] at (axis cs:1,1.4) {$e_1=1$};
  \node[right,font=\scriptsize] at (axis cs:3,3.9) {$e_3=2$};
  \node[right,font=\scriptsize] at (axis cs:5,5.4) {$e_5=1$};
\end{axis}
\end{tikzpicture}

\captionof{figure}{Data points, a model and the errors (red). The MSE is the average of the \emph{squares} of the red lengths.}\label{fig:b1_9}
\end{center}

\begin{aibox}
A neuron's weighted sum is exactly the dot product of Chapter~\ref{ch:dotproduct}. During training, a network minimises a loss such as the MSE: a sum over all training examples.
\end{aibox}

\section*{Practice problems}
\begin{enumerate}[label=\textbf{\thechapter.\arabic*}]
  \item Write out and compute $\displaystyle\sum_{i=1}^5 i^2$.
  \item Compute $\displaystyle\sum_{i=1}^4(2i+1)$.
  \item Write $3+6+9+12+15$ using sigma notation and compute the value.
  \item Compute the weighted sum $z=\sum w_ix_i$ for $w=(2,-1,0.5)$ and $x=(3,4,2)$. What is the result after adding the bias $b=-1$?
  \item Find the mean and the variance $\frac1n\sum(x_i-\bar x)^2$ of the numbers 4, 8, 6, 2 and 10.
  \item Compute the MSE for the true values $y=(1,0,1,1)$.\\ The predictions are $\hat y=(0.8,0.3,0.6,0.9)$.
  \item $\star$ Show that $\displaystyle\sum_{i=1}^n(x_i-\bar x)=0$ for any data.
  \item $\star$ Compute $\displaystyle\sum_{i=1}^{100}i$ by pairing first and last numbers ($1+100$, $2+99$, \dots). Check it with $\frac{n(n+1)}{2}$.
  \item \texttt{[code]} Write a Python loop that computes the MSE of Problem 3.6.
\end{enumerate}

% =====================================================
